\cotefile{46}
\scanpage{81}
% archive sheet number: 80
% block text | lines 1-2
Soit $\Gamma$ un groupe, $G$ un $\Gamma$-groupe, $M$ un $\Gamma$-\deleted{\unsure{groupe abélien}}\inserted{ensemble}. Supposons que $G$ opère sur $M$, et qu'on ait
% block display | line 3
\[
\unsure{\gamma|\sigma}(gm)=(\gamma g)(\gamma m)
\]
% block text | line 4
Soit alors $f$ un $1$-cocycle de $\Gamma$ à valeurs dans $G$
% block display | line 5
\[
f(\gamma\gamma')=f(\gamma)\cdot\gamma f(\gamma')
\]
% block text | lines 6-8
La donnée de $f$ permet de ``tordre'' $M$, de façon \unsure{précise} de définir une autre façon d'opérer de $\Gamma$ sur $M$, savoir
% block display | line 8
\[
\overset{\deleted{\unsure{\sim}}}{\gamma}_f m=f(\gamma)\cdot\gamma\cdot m
\]
% block display | lines 9-11
\[
\left[\begin{aligned}
\text{Car }\quad\gamma_f(\overset{\deleted{\unsure{\sim}}}{\gamma}'_f m)&=f(\gamma)\gamma(f(\gamma')\gamma'm)=f(\gamma)\gamma f(\gamma')\,\gamma\gamma'm\\
&=f(\gamma\gamma')\gamma\gamma'm=\overset{\deleted{\unsure{\sim}}}{(\gamma\gamma')}_f m\\
e_f m&=f(e)e\cdot m=m
\end{aligned}\right]
\]
% block text | lines 12-17
On désignera par $M_f$ le nouveau \deleted{\unsure{module}}\inserted{$\Gamma$-ensemble} ainsi construit; en fait, on ne s'est pas servi de la loi de groupe de $M$; en particulier, faisant opérer $G$ dans lui-même par automorphismes intérieurs, on \deleted{\unsure{trouve sur $G_f$ un module|a sur $G_f$ un module}} voit que
% block display | line 18
\[
\gamma(\sigma(g)\cdot g')=\sigma(\gamma g)\,\gamma g'
\]
% block text | lines 19-20
car le \deleted{\unsure{premier}}\inserted{premier} \unsure{mbr|mbre} est $\gamma(gg'g^{-1})=\gamma g\,\gamma g'(\gamma g)^{-1}$, c'est à dire $=$ le second membre. Ainsi, posant
% block display | line 21
\[
\gamma_f g=f(\gamma)\gamma g f(\gamma)^{-1}
\]
% block text | line 22
on obtient sur $G$ une nouvelle loi de $\Gamma$-groupe.

% block text | line 23
Notons \unsure{enfin} \unsure{.|:} nous avons
% block display | line 24
\[
\gamma_f(gm)=\gamma_f(g)\,\gamma_f(m)
\]
% block text | line 25
car le \deleted{deuxième}\inserted{premier} mbr est
% block display | lines 26-27
\[
\begin{aligned}
&f(\gamma)\gamma(gm)=f(\gamma)\gamma(g)\,\gamma(m)\\
&f(\gamma)\gamma g f(\gamma)^{-1}f(\gamma)\gamma(m)
\end{aligned}
\]
% block text | lines 26-27 right
et le deuxième \unsure{i.e.} pareil.
\scanpage{82}
% archive sheet number: 81
% block text | lines 1-2
Quel est \unsure{donc} $M_f^\Gamma$ ? C'est l'ensemble des $m\in M_f=M$ tels que
% block display | line 3
\[
f(\unsure{\gamma|\sigma})\gamma m=m\qquad\text{pour }\unsure{\mathrm{tt}|\forall}\ \gamma\qquad\text{i.e.}
\]
% block display | line 4
\[
\boxed{\gamma m=f(\gamma)^{-1}m\qquad\text{pour }\forall\ \gamma\in\Gamma}
\]
% block text | lines 5-6
On les appelle parfois éléments $f$-automorphes du $G$-$\Gamma$-ensemble $M$.

% block text | lines 7-13
Si $M$ est muni \inserted{en plus} de structures diverses invariantes par $\Gamma$ et $G$, alors il en est de $=$ de $M_f$ (qui conserve les structures qu'il y avait sur $M$). Si $M$ était p.ex. un $G$-groupe abélien, \deleted{\unsure{il est|et}} un $\Gamma$-groupe abélien, il en est de $=$ de $M_f$ etc. Dans ce cas, on peut s'intéresser, plus généralement, aux $H^i(\Gamma,M_f)$.

% block text | lines 14-17
Soit $f'$ \unsure{une|un} autre \unsure{$1$-cocycle|cocycle} de $\Gamma$ à valeurs dans $G$. Supposons $f'\sim f$, i.e. $f'\gamma=a^{-1}f(\gamma)\gamma a\quad(a\in G)$. On a alors un isomorphisme $u_a:M_{f'}\longrightarrow M_{f\deleted{\unsure{'}}}$, en \deleted{faisant correspondre}\inserted{posant}
% block display | line 18
\[
u_a m=a m
\]
% block text | line 19
Il faut vérifier
% block display | lines 20-22
\[
\begin{aligned}
u_a\gamma_{f'}m&=\gamma_f u_a m\\
\text{i.e.}\qquad a f'(\gamma)\gamma m&=f^{\deleted{\unsure{'}}}(\gamma)\gamma(am)\\
\text{i.e.}\qquad aa^{-1}f(\gamma)\gamma(a)\gamma(m)&=f(\gamma)\gamma(am)\ ;\ \text{c'est vrai.}
\end{aligned}
\]
% block text | lines 23-25
Si on veut que $\unsure{M_g|M_f}\simeq M_{f'}$ pour $\mathrm{tt}$ $G$-$\Gamma$-ensemble $M$, il \uline{faut} que $f\sim f'$ comme on voit en prenant $M=G$, $G$ y opérant par \unsure{translations gauches}.
\scanpage{83}
% archive sheet number: 82
% block text | top addition and lines 1-3
\inserted{Supposons que $M$ soit un groupe abélien.} La façon habituelle de construire des éléments de $M_f$ (du moins en supposant $\Gamma$ \uline{fini}) est \deleted{\unsure{par}} en prenant des moyennes
% block display | line 4
\[
N_f m=\sum_{\unsure{\gamma|\sigma}\in\Gamma}\gamma_f m=\sum_{\gamma\in\Gamma}f(\gamma)\gamma\cdot m
\]
% block text | lines 5-8
Supposons en particulier que $G$ soit l'ens des éléments inversibles d'un \uline{anneau} $A$, et que $M$ soit un $A$-module. Soit \deleted{\unsure{$N_f$ l'application de $A$|$N_f$ l'opérateur de $A$}}\inserted{pour \deleted{\unsure{$\forall$}} $u\in A$} \deleted{défini par}
% block display | line 8
\[
N_f u=\sum_{\gamma\in\Gamma}f(\gamma)\gamma(u)
\]
% block margin note | right of line 8
\margin{$\boxed{N_f(u)\deleted{\unsure{\in}}\in A_f^\Gamma}$}
% block text | line 9
Alors
% block display | line 9
\[
\deleted{\unsure{N_f u|N_f m}}xM^\Gamma\unsure{\subset|\in}M_f^\Gamma
\]
% block margin note | right of lines 9-10
\margin{En particulier, \unsure{$=$}}
% block display | line 10
\[
\text{i.e.}\qquad f(\gamma)\gamma(x)=\deleted{\unsure{f(\gamma)}}x\qquad\text{\unsure{où|si} }x=N_f u
\]
% block text | lines 11-14
[Si le groupe $\Gamma$ est infini, on essaie de faire des moyennes infinies, quand on arrive à leur donner un sens. C'est la \deleted{\unsure{méthode}} construction des fonctions \unsure{zeta-fuchsiennes|zêta-fuchsiennes|thêta-fuchsiennes} de Poincaré.]

% block text | small addition between lines 14 and 15
\inserted{Si on peut s'arranger pour trouver $u$ tel que $x$ soit inversible, on aura $f(\gamma)=x\gamma(x)^{-1}$ donc $f\sim0$}

% block text | line 15
\uline{Exemple} Supposons que $\Gamma$ \uline{\unsure{fini}} opère \uuline{fidèlement} dans \inserted{un} \uline{corps} $K$.

% block text | lines 16-17
Soit $A$ une algèbre de dim finie sur $K$, et supposons que $\Gamma$ opère dans $A$, de façon que
% block display | line 18
\[
\gamma(\lambda a)=(\gamma\cdot\lambda)(\gamma\cdot a)
\]
% block text | lines 19-23
Je dis qu'alors $H^1(\Gamma,G)\simeq(e)$, où $G$ est le groupe des éléments inversibles de $A$. D'après ce qui précède, il suffit de trouver un $\unsure{u|a}\in A$ tel que $\sum f(\gamma)\cdot\gamma u$ soit inversible, ce qui s'exprime par une \unsure{équation}.
\scanpage{84}
% archive sheet number: 83
% block display | line 1
\[
P\bigl(\sum f(\unsure{\gamma|\sigma})\cdot\gamma a\bigr)\ne0
\]
% block text | lines 2-9
où $P$ est un polynôme à coeff. dans $K$ \deleted{défini} défini sur l'espace $K$-vectoriel $A$, \unsure{p.ex. en des} coordonnées \deleted{\unsure{$a$}} du point courant, p.ex. à l'aide d'une base. On montre d'abord qu'on peut trouver une base formée d'éléments invariants par $\Gamma$ [c'est un lemme, qui n'utilise que la structure \deleted{\unsure{d'algèbre}}\inserted{$K$-$\Gamma$-espace} vectoriel de $A$, -- on n'utilise pas que la dimension soit finie]

% block text | line 10
Si $a=(a_i)_{1\leq i\leq\unsure{n}}$, on a donc
% block display | line 11 cancelled
\[
\deleted{P\bigl(\sum f(\gamma)\gamma a\bigr)=\unsure{Q|R}\bigl(f_1(\gamma),\ldots,f_n(\gamma)\bigr);}
\]
% block margin note | right of lines 10-11
% slanted
\margin{$R$ un polynôme en $n|\Gamma|$ variables}
% block display | line 12
\[
P\bigl(\sum f(\gamma)\gamma(a_\gamma)\bigr)=R\bigl(\{\{\gamma a_\gamma^1,\ldots,\gamma a_\gamma^n\}\}_{\gamma\in\Gamma}\bigr)
\]
% block text | lines 13-17
Supposons $K$ infini. Si l'\unsure{équation|expression} précédente était nulle \inserted{en y faisant $a_\gamma^i=a^i$ pour $\forall\gamma$ (et tout $i$)} quels que soient $a^1,\ldots,a^n\in K$, alors d'après le th. d'indépendance algébrique des automorphismes, alors $R$ est identiquement nul, en particulier on a \deleted{\unsure{$\sum$}} aussi
% block display | line 18
\[
\deleted{\unsure{R(\{\gamma a_\gamma^i}}\qquad P\bigl(\sum f(\gamma)\gamma(a_\gamma)\bigr)=0
\]
% block text | lines 19-20
quel que soit le choix des $a_\gamma$, d'où (faisant $a_{\gamma_0}=1$ et $\unsure{a_{\gamma_0}|a_\gamma}=0$ si $\gamma\ne\gamma_0$)
% block display | line 21
\[
P(f(\gamma_0))=0,\quad\text{absurde.}
\]
% block text | lines 21-23
/ Dans le cas où $K$ est fini, il faut utiliser \unsure{les résultats|le résultat} de Lang.....
\scanpage{85}
% archive sheet number: 84
% block text | line 1 heading
\deleted{\unsure{Corollaire}} \uline{Cas particulier}\qquad\uline{$H^1(\Gamma,K^*)=0$}.

% block text | lines 2-4
Ici, il suffit de vérifier qu'on peut trouver $x\in K^*$ tel que $\sum f(\unsure{\gamma|\sigma})\gamma x\ne0$, et l'indépendance \uline{linéaire} des automorphismes est suffisante.

% block text | lines 5-9
\uline{Généralisation} Soit $\underline{O}$ un anneau $\frac12$ local noethérien, et $\Gamma$ un groupe fini opérant \uline{fidèlement} sur $\underline{O}$ et \deleted{par} transitif sur les idéaux maximaux. Soit $\Gamma_d$ le \deleted{sous} groupe de stabilité d'un idéal maximal $\underline{m}$ de $O$, donc $\Gamma_d$ opère dans $O/m=k$.

% block text | lines 10-11
Soit $\Gamma_i$ le noyau de la représentation de $\Gamma_d$ \deleted{\unsure{dans le gr.|dans les}}\inserted{par} automorphismes du corps $O/m$.

% block text | lines 12-13
Soit enfin $A$ une $\underline{O}$-algèbre \deleted{\unsure{finie} \inserted{de type fini}} \deleted{Alors} où \inserted{$\Gamma$ opère tel que} \deleted{\unsure{$a\in A$}} \deleted{L'application naturelle} $\gamma(\lambda a)=(\gamma\lambda)(\gamma a)$. Alors les applications
% block diagram | line 14
\diagram{\begin{tikzcd}[ampersand replacement=\&,column sep=small]
H^1(\Gamma,A^*)\arrow[r]\&H^1(\Gamma_d,A_m^*)\arrow[r]\&H^1(\Gamma_i,A_m^*)\deleted{\unsure{/mA}}\arrow[r,"\deleted{\unsure{\longrightarrow}}" description,draw=none]\&\deleted{\unsure{H^1(\Gamma_i,(A/mA)^*)}}
\end{tikzcd}}
% block text | lines 15-19
sont \uline{injectives}. [N.B. \deleted{\unsure{$H^1(\Gamma_i,A/mA)$} \illegible{2}} dans \unsure{les cas les plus importants|le cas le plus important}, \deleted{\unsure{$A/mA$}} $\Gamma_i$ opère trivialement sur $A/mA$, donc $H^1(\Gamma_i,(A/mA)^*)$ est formé des classes d'homomorphismes -- à aut. intérieurs \inserted{de $A/mA$} près -- de $\Gamma_i$ dans $(A/mA)^*$]

% block text | lines 20-22
Soient $f,g\in Z^1(\Gamma,A)$, supposons qu'elles soient cohomologues dans \deleted{$A_m$}\inserted{$\unsure{H^1(\Gamma_d}$}, i.e. qu'il existe $a,\deleted{\unsure{b}}\in A$, $\lambda\in O\setminus m$, \deleted{\unsure{$b\in A$} \illegible{3}} $\unsure{b'|b}\in A$, $\mu\in O\setminus m$, tels que
% block display | line 23
\[
\deleted{(a\lambda^{-1})(b\mu^{-1})=}\quad\lambda^{-1}\mu^{-1}ab=\inserted{\lambda^{-1}\mu^{-1}ba=1}
\]
% block text | line 23 right and line 24 left
dans $A_m$ i.e. il existe $\nu\in O\setminus m$ tel que
% block display | line 24
\[
\deleted{\unsure{\nu\lambda\mu}}\nu ab=\nu\deleted{\unsure{\lambda\mu}}ba=\nu\lambda\mu\ ;
\]
% block text | line 24 right
et \unsure{les équations}
% block display | bottom insertion
\[
\inserted{g(\gamma)=(\unsure{a|a'}\lambda^{-1})^{-1}f(\gamma)\gamma(a\lambda^{-1})}
\]
% block text | line 25 right
dans $A_m$, i.e.
\scanpage{86}
% archive sheet number: 85
% block text | lines 1-2
Soient $\widehat{O},\widehat{A}$ les complétés de $O,A$, alors $\Gamma$ y opère de façon canonique. On a
% block display | line 3
\[
\widehat{O}=\bigoplus_i\widehat{O}_{m_i}\qquad\widehat{A}=\bigoplus_i\underbrace{\widehat{A}_{m_i}}_{B_i}
\]
% block text | lines 4-6
et $\Gamma_d$ \unsure{invarie} $B_0$, tandis que $\Gamma\bmod\unsure{\Gamma_i|\Gamma_d}$ permute les $B_i$ entre eux de façon simplement transitive. Il en résulte que
% block display | line 7
\[
H^1(\Gamma,\widehat{A}^*)\simeq H^1(\Gamma_d,\widehat{A}_m^*)
\]
% block text | lines 8-11
d'après un lemme bien connu (valable aussi en cohomologie non commutative). Il s'ensuit, puisque $f$ et $g$ définissent le $=$ élément de $H^1(\Gamma_d,\widehat{A}_m^*)$, \deleted{qu'il existe un $u\in A_m^*$} \inserted{[on regarde le diagramme commutatif}
% block text | left of diagram
\deleted{tel que}
% block diagram | lines 12-14
\diagram{\begin{tikzcd}[ampersand replacement=\&]
H^1(\Gamma,A^*)\arrow[r]\arrow[d]\&H^1(\Gamma_d,A_m^*)\arrow[d]\\
H^1(\Gamma,\widehat{A}^*)\arrow[r,"\sim"]\&H^1(\Gamma_d,\widehat{A}_m^*)
\end{tikzcd}}
% block text | lines 14-16
\inserted{]} qu'ils ont $=$ image dans $H^1(\Gamma,\widehat{A})$ donc qu'il existe un $u\in\widehat{A}^*$ tel qu'on ait
% block display | line 17
\[
g(\unsure{\gamma|\sigma})=u^{-1}f(\gamma)\gamma(u)
\]
% block text | lines 18-19 cancelled block
\deleted{Par \unsure{conséquent}, on peut, pour \unsure{$\forall$|toute} puissance $\unsure{\mathfrak r|\mathfrak m|\mathfrak b}^{N}$ du radical de $O$, trouver $v\in A$ tel que}
% block display | line 20 cancelled block
\[
\deleted{g(\gamma)\equiv u^{-1}f(\gamma)\gamma(u)\quad\bmod\mathfrak r^N A}
\]
% block text | line 21
Considérons l'application \inserted{$\unsure{O^\Gamma}$-}\deleted{\unsure{linéaire}}
% block display | line 22
\[
\varphi:A\longrightarrow A^\Gamma
\]
% block text | line 23
définie par
% block display | line 23
\[
\varphi(u)=\bigl(ug(\gamma)-f(\gamma)\gamma(u)\bigr)_{\gamma\in\Gamma}
\]
\scanpage{87}
% archive sheet number: 86
% block text | lines 1-5
Elle est linéaire \unsure{p.r.|par r.} à $O^\Gamma$. Or $O^\Gamma$ est un anneau $\frac12$ local noethérien, et $O$ est \inserted{module} de type fini sur $O^\Gamma$ (bien connu...), \unsure{enfin|alors} la topologie naturelle sur $A$ (donc $A^\Gamma$) est la $=$, qu'on le considère comme $O$-\inserted{module} ou comme $O^\Gamma$-module.

% block text | lines 6-10
Le noyau $N$ de l'application $\varphi$ a donc comme adhérence, dans $\widehat{A}$, le noyau $\widehat{N}$ de l'application prolongée par continuité. Or nous avons vu que ce dernier contient un élément \uline{inversible} $u$. \deleted{\unsure{Il en résulte aussitôt que}} \inserted{Soit $v\in N$ tel que}
% block margin note | left of lines 6-10
% slanted
\margin{\unsure{dégager} un lemme : si quel que soit $n$, on a $f\sim g\bmod m^n A$ \deleted{\unsure{démontrer|dégager}} alors $f\sim g$}
% block display | line 11
\[
v\equiv u\quad(\unsure{\mathfrak r|\mathfrak m|\mathfrak b}A)\qquad\text{où }\mathfrak r\text{ est le radical de }O.
\]
% block text | line 12
Alors $\unsure{u|v}$ est aussi inversible, et on a bien
% block display | line 13
\[
g\unsure{\gamma|\sigma}=u^{-1}f(\gamma)\gamma(u).
\]
% block margin note | left of lines 14-15
\margin{$\Vert$}
% block text | lines 14-15
Cela prouve que $H^1(\Gamma,A^*)\longrightarrow H^1(\Gamma_d,A_m^*)$ et $H^1(\Gamma,A^*)-H^1(\Gamma,\widehat{A}^*)$ sont injectifs.

% block text | lines 16-17
Supposons alors que $\underline{O}$ soit un \uline{anneau local} (donc $\Gamma=\Gamma_d$), \deleted{prouvons que} \inserted{\deleted{Considérons le diagramme}}
% block diagram | lines 18-20 cancelled block
\diagram{\begin{tikzcd}[ampersand replacement=\&]
\deleted{H^1(\Gamma,A^*)}\arrow[r]\arrow[d,"\deleted{\varphi}"]\&\deleted{H^1(\Gamma_i,A^*)}\arrow[d,"\deleted{\varphi_i}"]\\
\deleted{H^1(\Gamma,A/m^*)}\arrow[r,dash,"\deleted{\psi}"]\&\deleted{H^1(\Gamma_i,A/m^*)}
\end{tikzcd}}
% block text | left of diagram
\deleted{est injectif}
% block text | lines 21-22 cancelled block
\deleted{Supposons que nous prouvions que $\varphi$ et $\psi$ sont bijectifs, alors il en résulte (\unsure{l'appliq. à} $\Gamma_i$) que $\varphi_i$ l'est}
\scanpage{88}
% archive sheet number: 87
% block text | lines 1-2 cancelled block
\deleted{aussi, et par suite aussi la première flèche \deleted{est} horizontale.}

% block text | lines 3-4 cancelled block
\deleted{\deleted{\unsure{Supposons}} Prouvons donc que $\varphi$ est injective. Prouvons par récurrence sur $n\geq1$ que \deleted{\unsure{si}}}
% block display | lines 5-6 cancelled block
\[
\begin{aligned}
&\deleted{\text{si}\qquad f\sim g\quad\bmod\unsure{\mathfrak r|\mathfrak m|\mathfrak b}^n A}\\
&\deleted{\text{alors}\qquad f\sim g\quad\bmod\mathfrak r^{n+1}A}
\end{aligned}
\]
% block text | line 7 cancelled block
\deleted{Pour ceci considérons la suite exacte des $\Gamma$-groupes}
% block diagram | line 8 cancelled block
\diagram{\begin{tikzcd}[ampersand replacement=\&,column sep=small]
\deleted{e}\arrow[r]\&\deleted{N_n}\arrow[r]\&\deleted{(A/\mathfrak r^{n+1}A)^*}\arrow[r]\&\deleted{(A/\mathfrak r^n A)^*}\arrow[r]\&\deleted{0}
\end{tikzcd}}
% block text | lines 9-12 cancelled block
\deleted{Si $x,y\in A^*\deleted{\unsure{/\mathfrak r^{n+1}A}}$ leurs images dans $A/\mathfrak r^n A$ sont égales si et seulement si $x-y\in\mathfrak r^n A$ i.e. $xy^{-1}-1\in\mathfrak r^n A$ \deleted{et de $=$} i.e. $xy^{-1}\in1+\mathfrak r^n A$, et de $=$ pour l'\unsure{égalité|identité} \deleted{\unsure{suivant}} mod $\mathfrak r^{n+1}A$.}

% block text | lines 13-14 cancelled block
\deleted{Donc $N_n$ s'identifie au groupe additif $\mathfrak r^n A/\mathfrak r^{n+1}A$}
% block text | lines 14-15 restart
Je veux prouver que si $f,g$ sont tels que
% block display | line 16
\[
g(\unsure{\gamma|\sigma})=u^{-1}f(\gamma)\gamma(u)\qquad\gamma\in\Gamma_i\tag*{(*)}
\]
% block text | lines 17-18
alors $f\sim g$, \inserted{\unsure{i.e.} \unsure{$-$}} Nous pouvons d'abord faire opérer $\Gamma$ différemment sur $A$, par
% block display | line 19
\[
\gamma_f(x)=f(\gamma)\gamma\cdot x f(\gamma)^{-1}
\]
% block text | line 20 left
On a donc
% block display | line 20
\[
\gamma_f(\lambda x)=f(\gamma)\gamma(\lambda)\gamma(x)f(\gamma)^{-1}=\gamma(\lambda)\gamma_f(x)
\]
% block text | lines 21-22
donc $A_f$ est un $(O_f,\Gamma)$-module. \deleted{\unsure{Donc|Dans}} Cela étant, les équations (*) \unsure{s'écrit|s'écrivent}
% block display | line 23
\[
g(\gamma)f(\gamma)^{-1}=u^{-1}\gamma_f(u)\qquad\gamma\in\Gamma_i
\]
% block text | line 24
et on cherche $v$ tel qu'on ait
% block display | line 25
\[
g(\gamma)f(\gamma)^{-1}=v^{-1}\gamma_f(v)\qquad\text{pour }\text{\uuline{$\unsure{\mathrm{tt}|\forall}$}}\ \gamma\in\Gamma.
\]
\scanpage{89}
% archive sheet number: 88
% block text | line 1
Or on voit que
% block display | line 2
\[
h(\unsure{\gamma|\sigma})=g(\gamma)f(\gamma)^{-1}
\]
% block text | lines 3-8
est un $1$-cocycle de $\Gamma$ à valeurs dans $A_f$. On est ramené à prouver qu'il est \deleted{trivial}\inserted{$\sim0$}, sachant qu'il est \deleted{trivial}\inserted{$\sim0$} sur $\Gamma_i$. \unsure{Aut.} nous revenons donc au cas où l'un des cocycles $f,g$ est l'unité, soit $g$. Mais dans ce cas, on utilise la suite exacte
% block diagram | line 9
\diagram{\begin{tikzcd}[ampersand replacement=\&,column sep=small]
e\arrow[r,dash]\&H^1(\Gamma/\Gamma_i,G^{\Gamma_i})\arrow[r,dash]\&H^1(\Gamma,G)\arrow[r,dash]\&H^1(\Gamma_i,G)
\end{tikzcd}}
% block text | line 10
et on est ramené à prouver que
% block display | line 11
\[
H^1(\Gamma/\Gamma_i,G^{\Gamma_i})=\{e\}
\]
% block text | lines 12-20
Or $G^{\Gamma_i}=(A^{\Gamma_i})^*$. Soit $\Gamma'=\Gamma/\Gamma_i$, $A'=A^{\Gamma_i}$, \deleted{donc} $O'=O^{\Gamma_i}$. On sait que $O'$ est \inserted{local} noethérien, que son corps de restes est isomorphe à celui de $O$ donc que $\Gamma/\Gamma_i=\Gamma'$ y opère \uline{fidèlement}. De plus on sait que $O$ est un $O'$-module noethérien, d'où résulte qu'il en est de $=$ de $A$, donc de $A'$. $A'$ est une $O'$-algèbre, et $\Gamma'$ y opère de façon compatible avec les opérations sur $O'$. On est ramené à prouver
% block display | line 21
\[
H^1(\Gamma,A^*)=(e)
\]
% block text | line 22
\deleted{\unsure{si $O$}} dans le cas particulier où $\Gamma_i=\{e\}$,
\scanpage{90}
% archive sheet number: 89
% block text | lines 1-5
i.e. $\Gamma$ opère \uline{fidèlement} sur le corps des restes. On a déjà vu (comme conséquence du th. d'indépendance algébrique des automorphismes, \unsure{resp.} le lemme de Lang) que $H^1(\Gamma,A/mA^*)=\unsure{\{e)|(e)}$. Prouvons alors par récurrence sur $n\geq1$ que
% block display | line 6
\[
H^1(\Gamma,A/m^n A^*)=\{e\}
\]
% block text | lines 7-8
Pour ceci, on \deleted{suppose} le suppose fait pour $n\geq1$, on écrit la suite exacte
% block diagram | line 9
\diagram{\begin{tikzcd}[ampersand replacement=\&,column sep=small]
e\arrow[r,dash]\&N_n\arrow[r,dash]\&A/m^{n+1}A^*\arrow[r,dash]\&A/m^n A^*\arrow[r,dash]\&e
\end{tikzcd}}
% block text | lines 10-14
et on constate qu'en tant que $\Gamma$-groupe, $N_n$ est isomorphe à $m^{n+1}A/m^n A$, donc à un espace vectoriel sur $k=A/m$. \unsure{Donc|Or} il est (en tant que $\Gamma$-groupe) isomorphe à $k^N$, donc $H^1(\Gamma,k^N)=H^1(\Gamma,k)^N=0$ en vertu du \uline{lemme} $H^i(\Gamma,k)=0$ pour $i>0$.

% block text | line 15
La suite exacte
% block diagram | line 16
\diagram{\begin{tikzcd}[ampersand replacement=\&]
H^1(\Gamma,N_n)\arrow[r,dash]\&H^1(\Gamma,A/m^{n+1}A^*)\arrow[r,dash]\&H^1(\Gamma,A/m^n A)^*
\end{tikzcd}}
% block text | line 17
montre donc que l'on a
% block display | line 18
\[
H^1(\Gamma,A/m^{n+1}A^*)^{\deleted{\unsure{*}}}=\{e\}.
\]
\scanpage{91}
% archive sheet number: 90
% block text | lines 1-2
Or, reproduisant un raisonnement déjà fait, on en conclut que $H^1(\Gamma,A^*)=\{e\}$.

% block text | line 3
On est donc ramené à étudier
% block display | line 4
\[
H^1(\Gamma,A^*)
\]
% block text | lines 5-9
quand $O$ est un \uline{anneau local}, $\Gamma$ opère \uline{\unsure{trivialt|trivlt}} sur $O/m$, et $A$ est une $(O$-$\Gamma)$-algèbre. \uline{Supposons} \inserted{\uline{de plus que $\Gamma$ opère trivialement sur $A/mA$.}} \inserted{et soit} Soit $p$ la caractéristique de $k=O/m$, $\Lambda$ un sous-groupe de Sylow du type $p$ de $\Gamma$ ($\Lambda=\{e\}$ si la caractéristique est nulle).
% block margin note | left margin beside lines 5-17
% slanted
\margin{Nous supposerons plus précisément que $A^\Gamma\longrightarrow A/mA$ est \uuline{surjectif} [ce qui est p.ex. le cas si $A=M_n(\underline{O})$, $\Gamma$ opérant \unsure{trivialt|trivialement} sur $\underline{O}/m$, en effet, on sait qu'alors $\underline{O}^\Gamma\longrightarrow\underline{O}/m$ est surjectif\inserted{??}.....]}
% block text | line 10
Je dis qu'alors l'application
% block diagram | line 11
\diagram{\begin{tikzcd}[ampersand replacement=\&]
H^1(\Gamma,A^*)\arrow[r]\&H^1(\Gamma,A/mA^*)\times H^1(\Lambda,A^*)
\end{tikzcd}}
% block text | lines 12-15
est \deleted{bijective} injective (en particulier, si l'ordre de $G$ est premier à la car. $p$, en particulier si $k$ est de car. $0$, alors $H^1(\Gamma,A^*)\mathbin{-}H^1(\Gamma,A/mA^*)$ est injectif)

% block text | lines 16-22
Soient $f$ et $g$ \deleted{\unsure{enfin}} deux cocycles donnant les $=^{\mathrm{es}}$ images dans $H^1(\Gamma,A/mA^*)$ et $H^1(\Lambda,A^*)$. Pour montrer qu'ils sont $\sim$, on peut \unsure{encore|comme} supposer $g\unsure{(\gamma)|\gamma|(\sigma)}=\unsure{e|1}$ pour $\forall\gamma$. \deleted{\unsure{Soit le}} \inserted{Il suffit de} prouver alors \deleted{\unsure{$f\sim0$}} que $f$ est $\sim0$ dans \deleted{\unsure{$H^1(\Gamma,$|$H^1$},} $A/m^n A^*$ pour tout $n$. On \unsure{récurre} sur $n$, supposant que c'est fait pour un $n\geq1$.
\scanpage{92}
% archive sheet number: 91
% block text | lines 1-2
La suite exacte de coh. montre que $f$ est dans l'image de
% block display | line 3
\[
H^1(\Gamma,m^nA/m^{n+1}A)\longrightarrow H^1(\Gamma,A/m^{n+1}A^*)
\]
% block text | lines 4-5
i.e. est $\sim$ \inserted{dans $A/m^{n+1}A^*$} à un cocycle de la forme suivante \deleted{$\unsure{1+g\gamma|1+g}$}
% block display | line 6
\[
1+g(\unsure{\gamma|\sigma})
\]
% block text | line 7
où $g(\gamma)\in m^nA$ satisfaisant à
% block display | line 8
\[
g(\gamma\gamma')\equiv g(\gamma)+\gamma g(\gamma')\qquad(m^{n+1}A)
\]
% block text | lines 9-11
\deleted{Par hypothèse, on} On peut dans le \unsure{raisonnement} supposer $m^{n+1}=0$.

\deleted{Donc} L'hypothèse \deleted{\unsure{implique|indique}} implique qu'on a
% block display | line 12
\[
1+g(\gamma)=u^{-1}\gamma(u)\qquad\text{pour }\gamma\in\Lambda
\]
% block text | line 13
i.e.
% block display | line 13
\[
\gamma(u)-u=ug(\gamma)\deleted{\unsure{\in|\equiv|=}}\qquad\text{pour }\gamma\in\Lambda.\quad\deleted{\unsure{Comme}}
\]
% block insertion | above line 14
\inserted{\deleted{$mg(\gamma)=0$}}
% block text | lines 14-16, cancelled boxed block
\deleted{\unsure{Comme}}
% block display | line 14, cancelled boxed block
\[
\deleted{mg(\gamma)=0,\quad\text{i.e.}\quad ug(\gamma)=\unsure{\varepsilon|\epsilon}(u)g(\gamma)}
\]
% block text | line 15, cancelled boxed block
\deleted{i.e.}
% block display | line 15, cancelled boxed block
\[
\deleted{g(\gamma)=\gamma(\varepsilon(u)^{-1}u)-\varepsilon(u)^{-1}u}
\]
% block text | line 16, cancelled boxed block
\deleted{\unsure{on rappelle|en rappelant} que $\Gamma$ opère \uline{\unsure{trivialt|trivlt}} sur $O/m$ !)}

% block text | lines 17-23, cancelled block
\deleted{Cela signifie que $\gamma\mapsto ug(\gamma)$ est $\sim0$ dans $H^1(\Lambda,A)$ (additif !). \deleted{\unsure{Donc}} Or, $A$ est un $(\Gamma$-$O)$-module de longueur finie, on prouve par récurrence sur la longueur d'un tel module $M$ que l'application $H^1(\Gamma,M)\longrightarrow H^1(\Lambda,M)$ est \uline{injective}. \unsure{Or|Car}}

% block text | lines 24-25, doubly cancelled bottom lines
\deleted{\deleted{\unsure{Il suffit si} $M$ est \unsure{de longueur un|simple}, \unsure{ce qui est bien connu} : \illegible{1} par la \unsure{suite|raison} \illegible{1}}}

% block margin note | left margin, written bottom to top beside lines 1-24
% margin note written sideways
\margin{\deleted{Dans ce \inserted{dernier} cas, \deleted{\unsure{on a}} en effet $mM=0$ (sinon on aurait $mM=M$, d'où $M=0$ par Nakayama) donc $M$ est \deleted{\unsure{en}} en fait un vectoriel sur $k$ dans lequel $\Gamma$ opère par transformations \uline{$k$-linéaires}. On sait bien que $H^1(\Gamma,M)\longrightarrow H^1(\Lambda,M)$ est alors injectif.}}
\scanpage{93}
% archive sheet number: 92
% block text | line 1, cancelled inset
\deleted{\unsure{Caractéristique}. Si on a}
% block diagram | line 2, cancelled inset
\diagram{\begin{tikzcd}[ampersand replacement=\&]
\deleted{0}\arrow[r]\&\deleted{M'}\arrow[r,dash]\&\deleted{M}\arrow[r,dash]\&\deleted{M''}\arrow[r]\&\deleted{0}
\end{tikzcd}}
% block text | line 3, cancelled inset
\deleted{et si \unsure{on a} \illegible{2} \unsure{propriété|injectivité} \unsure{pour}}
% block text | lines 1-8, cancelled surrounding block
\deleted{\unsure{provient} du fait que, si $\unsure{\delta|d}=[\Gamma:\Lambda]$, la multiplication par $\delta$ est une bijection dans $A$ (car \unsure{ça l'est} dans les $m^iA/m^{i+1}A$). Par suite, on \deleted{\unsure{peut trouver}} a un $v\in A$ tel que}
% block display | line 9, cancelled surrounding block
\[
\deleted{\unsure{\gamma|\sigma}v-v=ug(\gamma)\qquad\text{pour }\underline{\forall}\gamma\in\Gamma}
\]
% block text | line 10
Soit $v\in A^\Gamma$ tel que $\unsure{\varepsilon|\epsilon}(v)=\varepsilon(u)$, on a donc
% block display | line 11
\[
g(\gamma)=\gamma(v^{-1}u)-v^{-1}u\qquad\text{d'où }g\sim0\text{ dans }H^1(\Lambda,A)
\]
% block text | lines 12-17
(additif !). Comme $\delta=[\Gamma:\Lambda]$ est tel que $x\mapsto\delta x$ dans $A$ est \uline{bijectif} (car il l'est dans les $m^iA/m^{i+1}A$, qui sont des espaces vectoriels sur $k$) \deleted{\unsure{on a}} $H^1(\Gamma,A)\mathbin{-}H^1(\Lambda,A)$ est \uline{injectif}, donc il existe $w\in A$ tel que
% block display | line 18
\[
g(\gamma)=\gamma(w)-w\qquad\text{pour }\underline{\forall}\gamma\in\Gamma.
\]
% block text | lines 18-22
Donc \deleted{$\unsure{1+g(\gamma)=wg|1+g(\gamma)=w\gamma}$} On peut (ajoutant au besoin \deleted{\unsure{un}} un élément de $A^\Gamma$) supposer $w\in A^*$, \deleted{\unsure{Soit donc}} \inserted{et} $=\varepsilon(w)=1$. Alors on a \deleted{\unsure{donc}} \unsure{on a}
% block display | lines 23-25
\[
\begin{aligned}
g(\gamma)&=\varepsilon(w)g(\gamma)=wg(\gamma),\qquad\text{d'où}\\
g(\gamma)&=w^{-1}\gamma(w)-1\qquad\text{d'où}\\
1+g(\gamma)&=w^{-1}\gamma(w)
\end{aligned}
\]
% block text | line 26
donc $1+g(\gamma)\sim0$ en $H^1(\Gamma,A^*)$.
\scanpage{94}
% archive sheet number: 93
% block text | upper-right title, lines 1-3
\begin{flushright}
Contre exemples à \unsure{passage} au \unsure{quotient}
\end{flushright}
\scanpage{95}
% archive sheet number: 94
% block margin note | boxed upper-right heading
\margin{Impossibilité \unsure{du|de} \unsure{passage} au quotient}
% block text | line 1
$G=GP(N)_{\unsure{\xi|S}}$ opérant sur $P^N_S$\qquad $S=\unsure{\operatorname{Spec}|\mathrm{Sp}}k$

% block text | lines 2-3
$V$ ouvert de $P^N_S$ sur lequel $G$ opère librement, supposons que $V/G$ existe ``effectivement''

% block text | lines 4-12
Alors on montre que le faisceau sur $V/G=W$ déduit de $\underline{L}=\underline{O}(+N+1)=(\underline{\Omega}^N_{P^N_S/S})^{-1}$ par ``descente'' est \uline{ample} sur $W$ [\unsure{utiliser|utilisant} le fait que $G$ n'a pas d'hom. non trivial dans $G_m$, et que \unsure{tt|tout} ouvert affine de $P^N_S$ est complément d'une hypersurface, \deleted{\unsure{donc}} définie par une section d'un $\underline{L}^{\otimes m}$ ($m>0$). À fortiori, $W$ est \uline{affine}. Donc si $G$ opère \inserted{librement} sur $V$ avec un \unsure{graphe} non fermé, $V/G$ n'existe pas.

% block text | lines 13-16, cancelled English block
\deleted{\deleted{\unsure{How to}} \inserted{so as} to free \unsure{arcaïc|archaïc} \unsure{K.-algV|K.-alg} from the \deleted{\unsure{heavy}} \unsure{weight|weights} of \inserted{\deleted{\unsure{unnecessary}} \deleted{\unsure{inadequate}} \unsure{inadequate}} concepts \deleted{conceptual background and technical \unsure{demands|details}} and \deleted{\unsure{new}} tools}
